% TEMPLATE-MILC-v3.tex - plantilla maestra para todas las guias MILC v3
% Ver scripts/generadores/build-guia-milc.py para la lista completa de
% claves esperadas. El script valida que no queden placeholders sin
% reemplazar antes de escribir el archivo final.

\documentclass[11pt,letterpaper]{article}
\usepackage[margin=1.75cm]{geometry}
\usepackage[table]{xcolor}
\usepackage{fontspec}
\usepackage[spanish]{babel}
\usepackage{tikz}
\usepackage[most]{tcolorbox}
\usepackage{tabularx}
\usepackage{array}
\usepackage{fancyhdr}
\usepackage{titlesec}
\usepackage{ragged2e}
\usepackage{enumitem}
\usepackage{microtype}

\setmainfont{Helvetica Neue}
\setsansfont{Helvetica Neue}
\setlength{\parindent}{0pt}
\setlength{\parskip}{6pt}
\hyphenpenalty=200
\exhyphenpenalty=200
\pretolerance=400
\tolerance=2000
\setlength{\emergencystretch}{1em}
\renewcommand{\arraystretch}{1.25}
\setlist{leftmargin=5mm,itemsep=1mm,topsep=1mm}
\newcolumntype{Y}{>{\RaggedRight\arraybackslash}X}

\definecolor{milcMagenta}{HTML}{E6007E}
\definecolor{milcVino}{HTML}{5A0038}
\definecolor{milcTurquesa}{HTML}{0093A5}
\definecolor{milcVerde}{HTML}{58A923}
\definecolor{milcMostaza}{HTML}{E5B400}
\definecolor{milcOcre}{HTML}{B86600}
\definecolor{milcNegro}{HTML}{241816}
\definecolor{milcGris}{HTML}{F7F7F4}
\definecolor{milcLinea}{HTML}{E8E2DD}
\definecolor{milcGradePrimary}{HTML}{0066FF}
\definecolor{milcGradeDark}{HTML}{003D99}
\definecolor{milcGradeSoft}{HTML}{D6E8FF}
\definecolor{milcGradeAccent}{HTML}{CFE7FF}
\definecolor{milcGradeText}{HTML}{FFFFFF}
\color{milcNegro}

\pagestyle{fancy}
\fancyhf{}
\lhead{\small Tecnología e Informática · Grado 11}
\rhead{\small Guía 27 · MILC v3}
\cfoot{\small\thepage}
\renewcommand{\headrulewidth}{0pt}

\newtcolorbox{titlebox}[1]{
  enhanced,colback=#1,colframe=#1,arc=7mm,boxrule=0pt,
  left=7mm,right=7mm,top=3mm,bottom=3mm,
  before skip=8pt,after skip=8pt,
  fontupper=\sffamily\bfseries\Large\color{white}
}

\newtcolorbox{softbox}[3]{
  enhanced,colback=#2,colframe=#1,borderline west={4pt}{0pt}{#1},
  arc=2mm,boxrule=.4pt,left=4mm,right=4mm,top=3mm,bottom=3mm,
  before skip=6pt,after skip=8pt,title={#3},coltitle=white,
  colbacktitle=#1,fonttitle=\sffamily\bfseries,fontupper=\RaggedRight,
  attach boxed title to top left={xshift=3mm,yshift=-2mm},
  boxed title style={arc=2mm, boxrule=0pt}
}

\newtcolorbox{drawbox}[2]{
  enhanced,colback=white,colframe=#1,arc=4mm,boxrule=1pt,
  left=5mm,right=5mm,top=4mm,bottom=4mm,before skip=8pt,after skip=8pt,
  title={#2},coltitle=white,colbacktitle=#1,fonttitle=\sffamily\bfseries,
  fontupper=\RaggedRight,
  attach boxed title to top left={xshift=4mm,yshift=-2mm},
  boxed title style={arc=3mm, boxrule=0pt}
}

\newtcolorbox{infoband}[1]{
  enhanced,colback=#1!8,colframe=#1!50,arc=2mm,boxrule=.5pt,
  left=4mm,right=4mm,top=2mm,bottom=2mm,before skip=4pt,after skip=4pt,
  fontupper=\small\RaggedRight
}

% Puente narrativo entre secciones (contrato editorial Conéctate).
% Da continuidad: "Acabas de X. Ahora vas a Y porque Z."
% Visualmente: banda izquierda en color del grado, texto en itálica pequeña.
\newtcolorbox{puentebox}{
  enhanced,colback=milcGradeSoft,colframe=milcGradeDark,
  borderline west={3pt}{0pt}{milcGradeDark},
  arc=0mm,boxrule=0pt,left=5mm,right=4mm,top=2.5mm,bottom=2.5mm,
  before skip=4pt,after skip=8pt,
  fontupper=\itshape\small\color{milcNegro}\RaggedRight
}
\newcommand{\puente}[1]{%
  \begin{puentebox}%
    \textcolor{milcGradeDark}{\textbf{\textrightarrow}}\hspace{2mm}#1%
  \end{puentebox}%
}

\newcommand{\linead}[1]{\noindent\color{milcLinea}\rule{#1}{0.5pt}\color{milcNegro}}
\newcommand{\respuesta}[1]{\par\vspace{2mm}\foreach \n in {1,...,#1}{\noindent\color{milcLinea}\rule{\linewidth}{0.5pt}\par\vspace{5mm}}\color{milcNegro}}
\newcommand{\checkbox}{\textcolor{milcGradeDark}{\fbox{\phantom{x}}}\hspace{5pt}}

\newcommand{\phasecard}[4]{
\begin{tcolorbox}[enhanced,colback=#3,colframe=#2,arc=3mm,boxrule=.5pt,height=3.05cm,valign=center,halign=center,left=1.5mm,right=1.5mm,top=2mm,bottom=2mm]
{\sffamily\bfseries\fontsize{22}{24}\selectfont\textcolor{#2}{#1}}\par
{\sffamily\bfseries\small #4}
\end{tcolorbox}
}

\begin{document}
\thispagestyle{empty}
\begin{tikzpicture}[remember picture,overlay]
  \fill[white] (current page.south west) rectangle (current page.north east);
  \path[fill=milcGradePrimary,rounded corners=16mm]
    ([xshift=.55cm,yshift=-.55cm]current page.north west) rectangle
    ([xshift=-.55cm,yshift=3.65cm]current page.south east);

  \node[anchor=north west,text=milcGradeText] at ([xshift=2.0cm,yshift=-1.85cm]current page.north west)
    {\sffamily\bfseries\fontsize{14}{16}\selectfont GRADO};
  \node[anchor=north west,text=milcGradeText,opacity=.82] at ([xshift=2.0cm,yshift=-2.75cm]current page.north west)
    {\sffamily\bfseries\fontsize{9}{11}\selectfont SERIE GUÍAS · TECNOLOGÍA E INFORMÁTICA};

  \begin{scope}[shift={([xshift=-3.05cm,yshift=-2.55cm]current page.north east)}]
    \fill[white,opacity=.80] (-.62,-.52) rectangle (.62,.52);
    \draw[milcGradeText,line width=1pt] (-.62,-.52) rectangle (.62,.52);
    \fill[milcGradeDark] (-.42,-.38) rectangle (-.22,.06);
    \fill[milcGradeAccent] (-.10,-.38) rectangle (.10,.24);
    \fill[milcGradePrimary!60!black] (.22,-.38) rectangle (.42,.38);
  \end{scope}

  \node[anchor=west,text=milcGradeText] at ([xshift=1.9cm,yshift=-12.85cm]current page.north west)
    {\sffamily\bfseries\fontsize{124}{124}\selectfont 11};
  \draw[milcGradeText,line width=1.7pt]
    ([xshift=4.52cm,yshift=-11.68cm]current page.north west) circle (.13cm);

  \node[anchor=west,text=milcGradeText,text width=8.7cm,align=left] at ([xshift=10.35cm,yshift=-11.6cm]current page.north west)
    {
     {\sffamily\bfseries\fontsize{19}{22}\selectfont Presupuesto ---\\cuánto siembras,\\cuánto cosechas}\\[4mm]
     {\sffamily\bfseries\fontsize{23}{26}\selectfont Guía 27-11-TIC}\\[2mm]
     {\sffamily\fontsize{13}{16}\selectfont Período 3 · Proyecto final emprendedor}\\[5mm]
     {\sffamily\bfseries\fontsize{11}{13}\selectfont Institución Educativa Sor María Juliana}\\[1mm]
     {\sffamily\fontsize{10}{12}\selectfont Autor: Dr. Álvaro Cárdenas Orozco}
    };

  \path[fill=milcGradeSoft,rounded corners=7mm]
    ([xshift=1.35cm,yshift=.85cm]current page.south west) rectangle
    ([xshift=-1.35cm,yshift=4.25cm]current page.south east);
  \draw[milcGradeDark,line width=.65pt,rounded corners=7mm]
    ([xshift=1.35cm,yshift=.85cm]current page.south west) rectangle
    ([xshift=-1.35cm,yshift=4.25cm]current page.south east);
  \node[anchor=center,text=milcNegro,text width=17.0cm,align=center] at ([yshift=2.55cm]current page.south)
    {\sffamily\fontsize{10}{12}\selectfont Metodología MILC · Apertura ancestral · Pensamiento computacional · Triángulo Dussel-Estoicismo-Floridi\\
    \textcolor{milcGradeDark}{\bfseries Producto:} Presupuesto detallado del MVP a 3 meses con (a) tabla de costos fijos y variables (con cifras realistas en pesos colombianos), (b) proyección de ingresos en 3 escenarios (pesimista, realista, optimista), (c) cálculo del punto de equilibrio mensual, (d) flujo de caja mes a mes de los próximos 3 meses, y (e) decisión por escrito sobre la viabilidad financiera mínima del proyecto};
\end{tikzpicture}
\mbox{}
\newpage

\section*{Apertura: Presupuesto --- cuánto siembras, cuánto cosechas}

\begin{softbox}{milcGradeDark}{milcGradeSoft}{Saber ancestral}
En las veredas del Valle del Cauca y en las fincas cafeteras del Quindío, ningún campesino sembraba sin antes hacer la \textbf{cuenta de la cosecha}. Esa cuenta no estaba escrita en libros: estaba en la cabeza del mayor y se transmitía a los hijos en las tardes de patio. Tenía 3 columnas claras: (1) \textbf{la siembra}, lo que se invertía antes de ver fruto (semilla, abono, jornales, herramienta, agua). (2) \textbf{la espera}, el tiempo que mediaba entre siembra y cosecha (3 meses para frijol, 9 meses para café, 1 año para plátano), durante el cual no entraba un peso. (3) \textbf{la cosecha}, lo que se esperaba recoger en kilos o en bultos, multiplicado por el precio del mercado más probable, descontando lo que se perdería por plaga, lluvia, transporte. El campesino que no hacía esta cuenta no era llamado \emph{emprendedor visionario}: era llamado \emph{loco} o \emph{novato}, y casi siempre perdía la finca al primer año adverso. La \textbf{phronesis económica campesina} sabía algo que el emprendimiento moderno suele olvidar: \textbf{el cálculo previo no enfría la pasión; la protege de la ruina}.
\end{softbox}

\begin{softbox}{milcTurquesa}{milcGris}{Saber contemporáneo}
El \textbf{presupuesto del emprendimiento digital} es exactamente la cuenta de la cosecha aplicada al MVP. Sus 3 piezas son las mismas con otros nombres: \textbf{costos} (lo que siembras), \textbf{horizonte temporal} (la espera), \textbf{ingresos proyectados} (lo que cosechas). A esto se suma una pieza moderna: el \textbf{punto de equilibrio} (break-even), el momento mensual en que los ingresos cubren los costos y dejan de exigir capital nuevo. La regla profesional: \textbf{ningún emprendimiento sin presupuesto a 3 meses sobrevive al cuarto mes}. Los costos se clasifican en dos familias: (a) \textbf{fijos}: ocurren cada mes existan o no ventas (alquiler de dominio, plataforma SaaS, sueldo propio si aplica, internet). (b) \textbf{variables}: dependen del volumen de operación (comisiones por transacción, costos de adquisición por usuario, materia prima, envíos). La phronesis financiera distingue ambos porque la viabilidad del proyecto depende de la \textbf{relación} entre costos fijos altos y margen variable bajo.
\end{softbox}

\begin{softbox}{milcMostaza}{milcGris}{Pregunta-puente}
\textit{¿Qué sabía el campesino al hacer la cuenta de la cosecha antes de sembrar, que el emprendedor novato olvida cuando se lanza con \emph{``ya veremos cómo nos va''}? ¿Y por qué un punto de equilibrio claro a 3 meses es señal de oficio profesional, no de pesimismo?}
\end{softbox}

\begin{softbox}{milcVerde}{milcGris}{Saber y saber hacer}
Al terminar podrás: (1) \textbf{identificar} los costos fijos y variables reales de tu MVP, con cifras en pesos colombianos verificables; (2) \textbf{explicar} la diferencia entre costos fijos, costos variables, ingresos, margen y punto de equilibrio aplicado a tu proyecto; (3) \textbf{aplicar} la fórmula del punto de equilibrio y construir el flujo de caja mes a mes de los próximos 3 meses; (4) \textbf{evaluar} con honestidad si tu proyecto es viable financieramente o necesita ajustar el modelo antes de seguir.
\end{softbox}

\small
\begin{tabularx}{\textwidth}{>{\bfseries}p{3.0cm}Y}
\rowcolor{milcGradePrimary!12}Autor & Dr. Álvaro Cárdenas Orozco\\
DBA & Diseño, implementación y sustentación de un proyecto de impacto local (MEN, T\&I)\\
Referentes & Dussel (analéctica · sur global) · Lean Startup · Floridi · Estoicismo\\
Producto final & Presupuesto detallado del MVP a 3 meses con (a) tabla de costos fijos y variables (con cifras realistas en pesos colombianos), (b) proyección de ingresos en 3 escenarios (pesimista, realista, optimista), (c) cálculo del punto de equilibrio mensual, (d) flujo de caja mes a mes de los próximos 3 meses, y (e) decisión por escrito sobre la viabilidad financiera mínima del proyecto\\
\end{tabularx}
\normalsize

\puente{Antes de empezar, mira el viaje completo. La sesión 5 te dio datos de uso del MVP y la sesión 6 los convirtió en pitch. Hoy le pones \textbf{números financieros} a la propuesta: cuánto cuesta sostenerla, cuánto puede generar, cuándo se equilibra. Las sesiones 8 (ética), 9 (versión final) y 10 (sustentación) se apoyarán en el presupuesto de hoy para hablar con seriedad de futuro.}

\begin{titlebox}{milcGradeDark}
Ruta MILC: así vamos a aprender
\end{titlebox}
\begin{tabularx}{\textwidth}{>{\centering\arraybackslash}X>{\centering\arraybackslash}X>{\centering\arraybackslash}X>{\centering\arraybackslash}X}
\phasecard{1}{milcVerde}{milcGris}{Escuta\\[-1mm]\small Observo, escucho y pregunto.} &
\phasecard{2}{milcTurquesa}{milcGris}{Sistematización\\[-1mm]\small Pensamiento computacional.} &
\phasecard{3}{milcMagenta}{milcGris}{Praxis\\[-1mm]\small Construyo y aplico.} &
\phasecard{4}{milcVino}{milcGris}{Evaluación\\[-1mm]\small Reflexión liberadora.}
\end{tabularx}


\puente{Antes de teorizar, vas a hacer un \emph{inventario honesto} de tus costos reales: qué herramientas usaste para el MVP, cuáles son gratuitas hasta cierto umbral, cuáles cobran desde el primer día, qué tiempo personal estás invirtiendo. Esa lista cruda es el insumo del presupuesto.}

\begin{titlebox}{milcVerde}
Fase 1 · Escuta
\end{titlebox}

\begin{softbox}{milcVerde}{milcGris}{Escena para escuchar}
\textbf{Actividad 1 · IDENTIFICA --- Inventario de costos reales del MVP} (20 min · individual). Toma tu MVP de la sesión 4 y lista en una tabla todos los costos que has tenido o tendrás en los próximos 3 meses. Reglas: (1) usa pesos colombianos (COP) con cifras verificables (no inventes, busca el precio real). (2) separa en \textbf{fijos} (cada mes, hayas vendido o no) y \textbf{variables} (dependen del volumen). (3) incluye tu \textbf{tiempo personal} valorado a un costo de oportunidad (mínimo 8.000 pesos por hora si calculas con SMMLV 2026). (4) cubre 7 categorías típicas: dominio/hosting, plataformas SaaS, marketing pagado, herramientas no-code (Tally, Zapier, ManyChat), papelería/impresión, comisiones de pago, tiempo personal. Si una categoría es cero pesos, ponla con 0 (no la omitas).
\end{softbox}

\checkbox Listé costos en pesos colombianos verificables.\\
\checkbox Separé fijos vs variables.\\
\checkbox Incluí el costo del tiempo personal.

\begin{infoband}{milcVerde}
\textbf{Lo que va al cuaderno (Actividad 1 · IDENTIFICA).} Título: «Actividad 1 --- Inventario de costos». \textbf{Formato:} tabla 7 filas × 4 columnas (Categoría~|~Fijo o variable~|~Costo mensual COP~|~Fuente del dato). Total mensual al final. \textbf{Sabes que terminaste cuando} los totales suman números reales que podrías defender en una sustentación.
\end{infoband}

\puente{Ya tienes los costos crudos. Ahora vas a entender la \textbf{anatomía financiera mínima}: las 3 columnas del presupuesto (costos, ingresos, flujo), la distinción entre fijos y variables, la fórmula del punto de equilibrio, los 3 escenarios de proyección (pesimista, realista, optimista) y los 6 errores típicos del presupuesto novato.}

\begin{titlebox}{milcTurquesa}
Fase 2 · Sistematización con pensamiento computacional
\end{titlebox}

\begin{softbox}{milcTurquesa}{milcGris}{Cuatro pilares aplicados al tema}
Una lista de costos sin estructura financiera es ruido contable. Para que tus números tomen forma profesional necesitas la \textbf{anatomía del presupuesto mínimo}: 3 piezas (costos, ingresos, flujo), 2 conceptos (margen, punto de equilibrio), 3 escenarios. Estas piezas son universales: las usa el campesino, el tendero y el emprendedor digital con las mismas matemáticas.
\end{softbox}

\small
\begin{tabularx}{\textwidth}{>{\bfseries\RaggedRight\arraybackslash}p{3.6cm}Y}
\rowcolor{milcTurquesa!90}\color{white}Pilar & \color{white}\bfseries Aplicación al tema\\
1. Descomposición & El presupuesto mínimo se descompone en 3 piezas: (1) \textbf{Costos totales}: suma de fijos + variables mensuales. (2) \textbf{Ingresos proyectados}: cuánto entra al mes según número de clientes multiplicado por ticket promedio. (3) \textbf{Flujo de caja}: cuánto queda neto cada mes (Ingresos menos Costos) y cómo se acumula mes a mes. Estas 3 piezas se calculan para los próximos 3 meses, no para 1 año (el primer año es demasiado especulativo en MVP).\\
2. Reconocimiento de patrones & El \textbf{punto de equilibrio} (break-even) tiene una fórmula sencilla: \emph{Punto de equilibrio (unidades) = Costos fijos dividido entre (Precio unitario menos Costo variable unitario)}. En palabras: ¿cuántas ventas necesito al mes para que los ingresos cubran los costos? Si tu modelo necesita 200 ventas mensuales para equilibrarse y tu MVP solo ha producido 5 inscripciones en 1 semana, tienes un problema financiero estructural. Detectarlo hoy es phronesis; descubrirlo en el mes 5 es ruina.\\
3. Abstracción & Los \textbf{3 escenarios de proyección} son disciplina profesional irrenunciable: (a) \textbf{Pesimista}: 30\% de la meta optimista (qué pasa si solo 3 de cada 10 cosas esperadas ocurren). (b) \textbf{Realista}: 60\% de la meta optimista (el escenario más probable según evidencia actual). (c) \textbf{Optimista}: 100\% de la meta optimista (si todo sale como esperas). La phronesis del emprendedor consiste en \textbf{tomar decisiones con base en el escenario realista, sostenibilidad en el pesimista, ambición en el optimista}. El novato vive solo en el optimista y se quiebra.\\
4. Algoritmo conceptual & Algoritmo: (1) calcula costos totales mensuales; (2) define precio unitario y costo variable unitario; (3) aplica fórmula del punto de equilibrio; (4) define meta optimista de ventas mensuales según evidencia del MVP; (5) calcula los 3 escenarios (30\%, 60\%, 100\%); (6) construye flujo de caja mes 1, 2, 3 para cada escenario; (7) declara la decisión: ¿proyecto financieramente viable o necesita ajustar modelo?\\
\end{tabularx}
\normalsize

\begin{drawbox}{milcTurquesa}{Anatomía del presupuesto a 3 meses}
\textbf{Costos fijos:} mensuales sin ventas. \\[1mm] \textbf{Costos variables:} unitarios por venta. \\[1mm] \textbf{Precio unitario:} cuánto cobras por unidad. \\[1mm] \textbf{Margen unitario:} Precio menos Costo variable. \\[1mm] \textbf{Punto de equilibrio:} Costos fijos divididos entre Margen. \\[1mm] \textbf{3 escenarios:} pesimista 30\% + realista 60\% + optimista 100\%. \\[1mm] \textbf{Flujo de caja:} Ingresos menos Costos mes a mes.
\end{drawbox}

\begin{softbox}{milcMostaza}{milcGris}{Errores comunes y su corrección}
(1) \textbf{Olvidar costos fijos invisibles}: dominio, suscripciones, tiempo personal. (2) \textbf{No separar fijos vs variables}: hace imposible calcular el punto de equilibrio. (3) \textbf{Proyectar solo escenario optimista}: el emprendimiento novato siempre proyecta el techo. (4) \textbf{Confundir ingresos con utilidad}: ingresos brutos no son utilidad neta. (5) \textbf{No incluir el costo del tiempo personal}: el emprendedor se autoexplota y cree que el proyecto es rentable. (6) \textbf{Punto de equilibrio nunca calculado}: imposible decidir si el modelo es viable sin esa cifra. (7) \textbf{Cifras inventadas sin fuente}: los números deben tener fuente verificable.
\end{softbox}

\begin{infoband}{milcTurquesa}
\textbf{Lo que va al cuaderno (Actividad 2 · EXPLICA).} Título: «Actividad 2 --- Anatomía financiera». \textbf{Formato:} (a) las 7 piezas de la anatomía aplicadas a tu MVP; (b) los 7 errores con marca del que más temes cometer; (c) cálculo escrito del punto de equilibrio con la fórmula y datos reales. \textbf{Sabes que terminaste cuando} podrías explicarle la lógica del punto de equilibrio a tu familia con ejemplo concreto.
\end{infoband}

\puente{Tienes el método. Toca aplicarlo: construir la tabla de costos, proyectar ingresos en 3 escenarios, calcular el punto de equilibrio, dibujar el flujo de caja mes a mes y declarar la decisión financiera por escrito.}

\begin{titlebox}{milcMagenta}
Fase 3 · Praxis con pensamiento computacional
\end{titlebox}

\begin{softbox}{milcMagenta}{milcGris}{Construyendo el producto con los cuatro pilares}
\textbf{Actividad 3 · APLICA + EVALÚA --- Presupuesto completo y decisión de viabilidad} (35 min · individual). Hora de aplicar. Construyes la tabla completa en Sheets o cuaderno: 3 meses, 3 escenarios, flujo de caja, punto de equilibrio. Cierras con la decisión documentada: \emph{proyecto viable, viable con ajustes, requiere pivote financiero, no viable}.
\end{softbox}

\small
\begin{tabularx}{\textwidth}{>{\bfseries\RaggedRight\arraybackslash}p{3.6cm}Y}
\rowcolor{milcMagenta!90}\color{white}Pilar & \color{white}\bfseries Aplicación al producto\\
Descomposición & Tu presupuesto se construye en 5 pasos: (1) tabla de costos fijos y variables mensuales (de la Actividad 1); (2) cálculo del punto de equilibrio en unidades vendidas; (3) proyección de ventas para 3 escenarios; (4) flujo de caja mes 1, 2, 3 para el escenario realista; (5) decisión de viabilidad financiera por escrito.\\
Patrones & Sigue el patrón \textbf{``escenario realista para decidir, pesimista para resistir''}: la decisión sobre seguir o ajustar se toma con base en el escenario realista (60\%); la pregunta sobre cuánto efectivo necesitas tener guardado se calcula con el pesimista (30\%). Esa disciplina dual evita dos errores opuestos: el optimismo ingenuo y el pesimismo paralizante.\\
Abstracción & Las 4 decisiones financieras posibles tienen criterios claros: (a) \textbf{Viable}: el escenario realista alcanza el punto de equilibrio en el mes 3 o antes. (b) \textbf{Viable con ajustes}: el escenario realista está cerca del punto de equilibrio pero requiere reducir 1-2 costos o subir el precio. (c) \textbf{Requiere pivote financiero}: el escenario realista no alcanza el punto de equilibrio en 6 meses; hay que cambiar segmento, precio o modelo. (d) \textbf{No viable como está}: incluso el optimista no alcanza el punto de equilibrio; el modelo necesita repensarse desde cero.\\
Algoritmo de armado & Algoritmo de cierre: (1) declara la decisión en 1 frase; (2) justifica con 2 cifras concretas (\emph{``decido que es viable con ajustes porque el punto de equilibrio son 80 ventas/mes y el escenario realista proyecta 60''}); (3) anota el cambio que harías si decidiste ajustar; (4) calcula cuánto efectivo necesitas para sostener el escenario pesimista los primeros 3 meses.\\
\end{tabularx}
\normalsize

\begin{drawbox}{milcMagenta}{Checklist antes de cerrar el presupuesto}
\checkbox Tabla de costos fijos y variables con cifras verificables. \\[1mm] \checkbox Costo del tiempo personal incluido. \\[1mm] \checkbox Punto de equilibrio calculado con fórmula. \\[1mm] \checkbox 3 escenarios proyectados (pesimista, realista, optimista). \\[1mm] \checkbox Flujo de caja mes 1, 2, 3 para escenario realista. \\[1mm] \checkbox Decisión de viabilidad escrita y justificada con cifras. \\[1mm] \checkbox Efectivo mínimo de respaldo calculado.
\end{drawbox}

\begin{softbox}{milcMagenta}{milcGris}{Plantilla de trabajo}
\textbf{Costos fijos.} \textit{Total mensual COP.} \\[1mm] \textbf{Costos variables.} \textit{Por unidad COP.} \\[1mm] \textbf{Precio.} \textit{Unitario COP.} \\[1mm] \textbf{Punto de equilibrio.} \textit{Costos fijos divididos entre margen unitario = unidades por mes.} \\[1mm] \textbf{Escenarios.} \textit{Pesimista 30\%, realista 60\%, optimista 100\%.} \\[1mm] \textbf{Flujo de caja.} \textit{Mes 1, 2, 3 en escenario realista.} \\[1mm] \textbf{Decisión.} \textit{Viable / Viable con ajustes / Pivote / No viable + justificación.}
\end{softbox}

\begin{infoband}{milcMagenta}
\textbf{Lo que va al cuaderno (Actividad 3 · APLICA + EVALÚA).} Título: «Actividad 3 --- Mi presupuesto a 3 meses». \textbf{Formato:} (a) tabla completa de costos, ingresos, flujo; (b) cálculo del punto de equilibrio; (c) los 3 escenarios; (d) decisión de viabilidad en 3 líneas. \textbf{Extensión:} 2 páginas. \textbf{Sabes que terminaste cuando} podrías presentarle el presupuesto a tu familia y defender por qué tu proyecto se sostiene o no.
\end{infoband}

\puente{Lo que sigue resume \emph{qué} entregas hoy: presupuesto a 3 meses + 3 escenarios + punto de equilibrio + flujo de caja + decisión de viabilidad. El criterio principal: que un familiar emprendedor o un asesor financiero, mirando tu hoja, pueda decir en 3 minutos \emph{``este proyecto se sostiene''} o \emph{``a este proyecto le faltan X y Y''}.}

\begin{titlebox}{milcMostaza}
Producto de la sesión
\end{titlebox}
\textbf{Presupuesto a 3 meses + punto de equilibrio + 3 escenarios + decisión documentada}.
\begin{softbox}{milcMostaza}{milcGris}{Criterios de calidad}
(1) Costos fijos y variables con cifras verificables. (2) Tiempo personal incluido en costos. (3) Punto de equilibrio calculado con fórmula. (4) 3 escenarios proyectados con honestidad. (5) Flujo de caja mes a mes para escenario realista. (6) Decisión de viabilidad escrita con justificación numérica.
\end{softbox}

\puente{Antes de cerrar, mira el presupuesto desde las cinco dimensiones humanas. Hacer la cuenta con rigor es respeto por uno mismo, por el inversionista futuro y por los usuarios; saltarse la cuenta es invitar a la ruina con buenos modales.}

\begin{titlebox}{milcVino}
Fase 4 · Evaluación liberadora — cinco dimensiones
\end{titlebox}

\begin{softbox}{milcVino}{milcGris}{Sentido de la evaluación}
Evaluar no es solo revisar una nota. Es reconocer qué aprendí, qué cambió en mí, qué emociones manejé, cómo me relaciono con la ciudad y la comunidad, y qué le dejaré a quien viene detrás.
\end{softbox}

\textbf{Mi reflexión liberadora en cinco dimensiones.} Completa cada frase iniciada:

\small
\begin{tabularx}{\textwidth}{>{\bfseries\RaggedRight\arraybackslash}p{3.4cm}Y>{\RaggedRight\arraybackslash}p{4.6cm}}
\rowcolor{milcVino}\color{white}Dimensión & \color{white}\bfseries Mi respuesta (frame starter) & \color{white}\bfseries Algo que descubrí\\
1. Personal & Lo que más cambió en mí fue \linead{6.5cm} & \linead{4.2cm}\\
2. Emocional & La emoción más fuerte fue \linead{2.5cm} y la regulé \linead{3cm} & \linead{4.2cm}\\
3. Ciudadana & En democracia esto importa porque \linead{6cm} & \linead{4.2cm}\\
4. Local (Cartago) & En mi barrio sería útil para \linead{6.5cm} & \linead{4.2cm}\\
5. Intergeneracional & A mi abuela/abuelo le diría \linead{6.5cm} & \linead{4.2cm}\\
\end{tabularx}
\normalsize

\begin{infoband}{milcVino}
\textbf{Para la próxima sustentación.} Vuelve a esta tabla en seis meses. Verás que las respuestas cambian — no porque lo aprendido cambie, sino porque tú cambias. La evaluación liberadora es una conversación contigo a través del tiempo.
\end{infoband}


\puente{Y para terminar, tres voces sobre la prudencia económica. Dussel sobre los modelos que se sostienen sin extraer, Séneca sobre la sobriedad del cálculo, Floridi sobre la transparencia financiera como ética profesional contemporánea.}

\begin{titlebox}{milcGradeDark}
Triángulo de pensamiento · cierre filosófico
\end{titlebox}

\begin{softbox}{milcVino}{milcGris}{Dussel · lente del nosotros}
\textit{``Un presupuesto que se sostiene sin extraer es ya una toma de posición ética.''}

Tu presupuesto puede sostenerse de dos maneras: (a) cobrando un margen justo al usuario por valor real entregado, o (b) extrayendo valor más allá del beneficio que devuelves (comisiones ocultas, dependencia generada, precios que aprovechan falta de opciones). La filosofía de la liberación pide la primera vía. Un presupuesto extractivo se ve rentable a corto plazo y se derrumba cuando el cliente encuentra opciones; un presupuesto justo se sostiene a largo plazo aunque el margen sea menor.

\textbf{¿Mi presupuesto cobra por valor entregado o se sostiene extrayendo más de lo que aporta? ¿Mi cliente, si viera todos mis números, los aprobaría?}
\end{softbox}
\linead{\linewidth}\\[1mm]
\linead{\linewidth}

\begin{softbox}{milcMostaza}{milcGris}{Estoicismo · Séneca · cuidado interior}
\textit{``Calcular antes es virtud; lamentar después es vicio del que se ahorra el cálculo previo.''}

Es tentador saltarse el presupuesto porque \emph{``ya veremos''}. La sobriedad estoica recomienda lo opuesto: \emph{calcula antes, lamenta menos}. El presupuesto no enfría la pasión: la protege. Quien se prohíbe el cálculo previo no está siendo audaz; está siendo perezoso. La phronesis financiera honra a quien hace la cuenta aunque duela, no a quien la evita para mantener el optimismo intacto.

\textbf{¿Estoy haciendo el cálculo aunque revele problemas, o lo estoy evitando para no enfrentar lo que muestra?}
\end{softbox}
\linead{\linewidth}\\[1mm]
\linead{\linewidth}

\begin{softbox}{milcTurquesa}{milcGris}{Floridi · lente de la infoesfera}
\textit{``La transparencia financiera es la nueva ética profesional en el oficio digital.''}

Un presupuesto que el cliente o el inversionista nunca podría ver sin escándalo no es modelo sostenible: es modelo opaco. La ética digital contemporánea exige que el emprendedor pueda mostrar sus cifras y defenderlas. Esa transparencia financiera te entrena para una práctica profesional que rige toda carrera técnica del siglo XXI: \textbf{el dato económico bien presentado y bien defendido}.

\textbf{Si mi familia, mi cliente o un inversionista vieran mi presupuesto completo, ¿lo aprobarían o detectarían algo que oculté?}
\end{softbox}
\linead{\linewidth}\\[1mm]
\linead{\linewidth}

\begin{titlebox}{milcVino}
Compromiso final
\end{titlebox}

\small
\begin{tabularx}{\textwidth}{>{\RaggedRight\arraybackslash}p{3.0cm}YYY}
\rowcolor{milcVino}\color{white}\bfseries Criterio & \color{white}\bfseries Lo logré & \color{white}\bfseries En proceso & \color{white}\bfseries Necesito apoyo\\
Comprensión del tema & Explico ideas centrales con mis palabras. & Reconozco ideas pero falta precisión. & Necesito repasar conceptos base.\\
Pensamiento computacional & Apliqué los 4 pilares al diseñar mi producto. & Apliqué algunos pilares. & Necesito releer la fase 2.\\
Aplicación y producto & Mi producto cumple los criterios y se puede revisar. & Producto funcional con ajustes pendientes. & Aún en construcción.\\
Reflexión 5 dimensiones & Respondí cada dimensión con honestidad. & Respondí algunas con detalle. & Mis respuestas son muy generales.\\
Triángulo de pensamiento & Conecté las 3 voces con mi trabajo real. & Conecté al menos una. & No relaciono las voces con mi proceso.\\
\end{tabularx}
\normalsize

\textbf{Hoy entendí que:}\\[1mm]
\linead{\linewidth}\\[1mm]
\linead{\linewidth}

\textbf{Mi compromiso para la próxima sesión es:}\\[1mm]
\linead{\linewidth}\\[1mm]
\linead{\linewidth}

\begin{softbox}{milcMostaza}{milcGris}{Cita de despedida · Marco Aurelio}
\textit{``El obstáculo en el camino se vuelve el camino. Lo que estorba al camino se vuelve camino.''}\\[1mm]
Cada error de hoy, bien observado, se convierte en pieza del camino que pisarás mañana.
\end{softbox}

\small
\begin{tabularx}{\textwidth}{>{\bfseries}p{3.6cm}Y}
\rowcolor{milcGradeDark!12}Nombre completo: & \linead{10cm}\\
Fecha: & \linead{4cm} \quad Curso/grupo: \linead{4cm}\\
Mi firma: & \linead{9cm}\\
\end{tabularx}
\normalsize

\begin{infoband}{milcGradeDark}
\textbf{Para la próxima sesión.} Llega con el presupuesto cerrado y la decisión documentada. En la sesión 8 vas a trabajar la \textbf{dimensión ética y legal} del proyecto: política de privacidad, términos de uso, responsabilidad sobre datos. Sin presupuesto sólido hoy, los compromisos éticos de mañana flotarán sin sostén económico.
\end{infoband}

\end{document}
